% 爱德华·康登（综述）
% license CCBYSA3
% type Wiki

本文根据 CC-BY-SA 协议转载翻译自维基百科\href{https://en.wikipedia.org/wiki/Edward_Condon}{相关文章}
\begin{figure}[ht]
\centering
\includegraphics[width=6cm]{./figures/748cfbd4f06fd931.png}
\caption{} \label{fig_ADHkd_1}
\end{figure}
爱德华·尤勒·康登（Edward Uhler Condon，1902年3月2日－1974年3月26日）是一位美国核物理学家、量子力学的先驱之一，并在第二次世界大战期间参与了雷达的研制工作，以及在曼哈顿计划中短暂参与了核武器的开发。以他命名的“弗兰克–康登原理”和“斯莱特–康登规则”均纪念了他在该领域的贡献。[1][2][3]

他曾于1945年至1951年间担任美国国家标准局（现为美国国家标准与技术研究院，NIST）的第四任局长。1946年，康登出任美国物理学会会长；1953年，他担任美国科学促进会会长。

在麦卡锡主义时期，康登是最早成为美国众议院非美活动调查委员会攻击目标的知名科学家之一。1948年，他被公开指控为“美国原子安全体系中最薄弱的一环”，理由是他掌握大量机密信息、与原子弹研制有关联，并被指涉嫌同情共产主义和苏联。康登事件成为当时反对麦卡锡主义者（尤其是科学界）关注的标志性案件之一，舆论广泛关注，他也得到了包括美国总统哈里·杜鲁门在内的众多知名科学家的公开支持。[4]

1968年，康登作为《康登报告》的主要作者再次广为人知。该报告是由美国空军资助的官方评估，得出的结论是“不明飞行物都有平凡的解释”。月球上的“康登环形山”即以他的名字命名，以纪念他的贡献。
\subsection{背景}
